% EGA V, English translation by P. Blass, 1987 (cote 24): archive cote 24, site page 44 (PDF page 45).
% typescript. Draft compared: none.
% AI-transcribed by Claude Opus 5.5 (claude-opus-5-5, adaptive thinking, effort high), Message Batches API, batch ega24b (msgbatch_01RGJJeHdR96U16ZWFuUjtDx), 2026-11-13; not yet checked by eye.
\scanpage{24/44}
\mspage{44}

x.  The announced characterization of singular zeros [illegible
ask AG] the elements of a smooth subscheme $\underline{P}(\hand{\unsure{\nu}}\ X/\underline{P})$ of $X_{\underline{P}}{}^{v}$ gives
immediately the following statement which deserves to appear as a
preliminary proposition if S = Spec k and if H is a hyperplane of
$\underline{P}$ then $Y = X \ \mathrm{x}_{\underline{P}}\ H$ is smooth over k of relative dimension (d - 1)
at the point $x \in \underline{Y}(k)$ (i.e. x is a non-singular zero i.e.
geometrically non-singular of the section {\o} of $O_X(1)$ defined by H)
if and only if H does not contain the image by {\o} of the tangent
space to X at x (relative to k) or as we say once more (if
$f{:}X \to \underline{P}$ is an immersion which allows us to identify X to a
subscheme of $\underline{P}$) if and only if H is not tangent to X at x.  This
follows trivially from the Jacobian criterion of smoothness or
from the definition of a singular zero, once we make precise the
sense of the statement, that is to say that we make precise how a
vector subspace of the tangent space to $\underline{P}$ at a point a(= f(x))
defines a linear subspace of $\underline{P}$ (in such a way that it makes sense
to say that H does not contain the soul vector** subspace): of
course this comes from the exact sequence (b) above which allows
to define a one-to-one correspondence between the set of vector
subspaces of the tangent space at a and the set of linear
subspaces to $\underline{P}$ (?)* [illegible ask AG] containing a.  This
correspondence anyway reduces to associating to a linear
subvariety passing through a its tangent space at a considered as
a subspace of the tangent space to $\underline{P}$ at a.

Such ``sorites'' grouped together with various ``sorites'' about
linear subvarieties and about grossmanians ought to be given in
one or to preliminary numbers or paragraphs of course announcing

\margin{** Check this Blass}
